\documentclass[11pt]{article}
\usepackage{microtype,amsmath,amssymb,amsthm,mathtools,booktabs,array}
\usepackage[margin=1.05in]{geometry}
\usepackage[colorlinks=true,linkcolor=blue!45!black,citecolor=blue!45!black,urlcolor=blue!45!black]{hyperref}
\usepackage{xcolor}
\hypersetup{pdftitle={The dimension-degree boundary for Gaussian moment vanishing},pdfsubject={Working mathematical manuscript}}
\newtheorem{theorem}{Theorem}[section]
\newtheorem{proposition}[theorem]{Proposition}
\newtheorem{lemma}[theorem]{Lemma}
\newtheorem{corollary}[theorem]{Corollary}
\theoremstyle{definition}\newtheorem{definition}[theorem]{Definition}
\theoremstyle{remark}\newtheorem{remark}[theorem]{Remark}
\newcommand{\E}{\mathbb E}
\newcommand{\C}{\mathbb C}
\newcommand{\tr}{\operatorname{tr}}
\newcommand{\wt}{\operatorname{wt}}
\newcommand{\G}{\mathsf G}
\DeclareMathOperator{\Span}{span}
\title{Stationary actions and the finite-critical ordinary-node exclusion}
\author{}\date{17 November 2026: internal written proof}
\begin{document}\maketitle
This proof covers the finite affine critical-set part of three ternary
cubic families,with all Gaussian means and cross terms retained.
It replaces separate escape-scale computations by an exact action
identity. The global rank input is classical:Pham's tame Milnor
decomposition and Roucairol's nonramified formal direct-image rank
formula. No novelty,independent replay,formal verification or external
review is claimed. Positive-dimensional critical sets remain outside
this proof and are listed exhaustively at its end.

For independent standard Gaussian variables define the formal series
\[
 \Phi_A(t)=\sum_{m\ge0}\mathbb E[AP^m]t^m/m!.
\]
Purity means $\Phi_1=1$. Gaussian integration by parts gives a horizontal
functional on the generic top twisted quotient
\[
 M=\mathbb C(t)[x_1,x_2,x_3]/
       \sum_i(\partial_i+tP_i-q_i)\mathbb C(t)[x_1,x_2,x_3],
 \qquad\nabla_t=\partial_t+P.
\]
Under an invertible affine coordinate change,the transformed Gaussian
functional and quadratic $q$ retain these identities;this is the
meaning of the full Gaussian pair below. Only rational seed restriction
and formal completion of this ambient module are used. No convergence
of the moment series,actual period choice or image-kernel computation
is assumed. Write $P=F_3+F_2+F_1+F_0$ for its homogeneous pieces.
\section{Algebraic reductions}
\begin{lemma}[The algebraic frame test]\label{lem:frame}
There are Gaussian coordinates $(Z,W,T)$ in which $P$ is affine linear
in $W$ if and only if there is a nonzero vector $v\in\C^3$ such that
\begin{equation}\label{eq:frame}
 v^Tv=0,\qquad \nabla F_3(v)=0,\qquad F_2(v)=0.
\end{equation}
\end{lemma}
\begin{proof}
Polarization gives
\[
 D_v^2P(x)=2\nabla F_3(v)^Tx+2F_2(v).
\]
The last two conditions in \eqref{eq:frame} say exactly that this
polynomial vanishes. Complete the nonzero null vector $v$ to a Witt
basis: choose a null partner with pairing one, and a unit vector in
their orthogonal complement. In these coordinates $D_v$ is the $W$
derivative. Hence $D_v^2P=0$ says that no second or higher power of
$W$ occurs. The converse reads off $v$ from the $W$ direction.
\end{proof}

\begin{lemma}[The ordinary-node test]\label{lem:nodecoercivity}
Let $P=F_3+F_2+F_1+F_0$ in $n\geq2$ variables. Suppose its
projective singular set is finite and every singular point $[v]$
has $\ker\operatorname{Hess}F_3(v)=\C v$. If $F_2(v)\ne0$ at
every such point, then
$\|\nabla P(x)\|\geq c\|x\|$ outside a compact set. 
\end{lemma}
\begin{proof}
Otherwise choose an unbounded sequence with
$\|\nabla P(x)\|/\|x\|\longrightarrow0$. A projective limit $[v]$
is singular for $F_3$. Choose a fixed complement to $\C v$ and write
$x=sv+y$, where $y=o(s)$. With $H=\operatorname{Hess}F_3(v)$,
homogeneity gives the exact expansion
\[
 \nabla P(sv+y)=sHy+\nabla F_3(y)+s\nabla F_2(v)
                              +\nabla F_2(y)+\nabla F_1.
\]
The map $H$ is injective on the chosen complement. Divide by $s$.
The error from $\nabla F_3(y)/s$ is $o(\|y\|)$, so the equation
first forces $y$ to stay bounded. A limit $y_0$ satisfies
$Hy_0+\nabla F_2(v)=0$. Pair bilinearly with $v$. Symmetry and
$Hv=0$ kill the first term, whereas Euler's identity gives
$v\cdot\nabla F_2(v)=2F_2(v)$, a contradiction.
\end{proof}
\begin{lemma}\label{lem:ordinarynormal}
Let $[w]$ be an ordinary node of a ternary $F_3$, with $F_2(w)=0$
and $w^Tw\ne0$. An invertible affine change of coordinates puts
the pair consisting of $P$ and the Gaussian quadratic into
\begin{align*}
 P={}&T(uv+h)+a u^3+d v^3+b u^2+c v^2+e u+f v+k,\\
 q={}&\tfrac12(Au^2+Bv^2+mT^2+2Cuv+2DuT+2EvT)
                  +\ell_u u+\ell_v v+\ell_T T+q_0,
\end{align*}
where $m\ne0$ and $\det\operatorname{Hess}q\ne0$.
All Gaussian means and cross terms must be retained. Its leading
curve is irreducible nodal when $ad\ne0$, a transverse line and
smooth conic when exactly one of $a,d$ vanishes, and three
nonconcurrent lines when $a=d=0$.
\end{lemma}
\begin{proof}
Take $T$ in the direction $w$ and choose a complementary plane.
The singular-point equations make $F_3$ linear in $T$; its
coefficient is a nondegenerate binary quadratic because the node
is ordinary. Factor that quadratic and choose its two linear
factors as coordinates $u,v$. Thus $F_3=Tuv+B_3(u,v)$.
Replace $T$ by $T$ plus a linear form in $u,v$ to remove the
$u^2v$ and $uv^2$ terms of $B_3$.

There is no $T^2$ term in $F_2$, precisely because $F_2(w)=0$.
The full coefficient of $T$ is therefore
$uv+\alpha u+\beta v+h_0$. Translations of $u,v$ make it $uv+h$.
The cubic terms outside $T$ remain pure cubes. A constant shift
of $T$ removes the remaining quadratic $uv$ term. This gives the
displayed $P$. The direction of $T$ has not changed, so $m\ne0$;
invertibility of the affine change preserves the nondegenerate
Gaussian Hessian, while translations and the shifts of $T$ produce
the means and cross terms shown.

If $ad\ne0$, the polynomial $Tuv+a u^3+d v^3$ is irreducible:
an independent factor of a polynomial linear in $T$ would divide
both $uv$ and $a u^3+d v^3$, whose gcd is one. Its gradient shows
that its sole singular point is the stated ordinary node. If
$a=0,d\ne0$, it is $v(Tu+d v^2)$, with a smooth conic meeting
the line transversely at two points; the other one-zero case is
symmetric. If both vanish, it is the coordinate triangle $Tuv$.
\end{proof}

\section{Finite-critical exclusion}
\subsection{Ordinary nodes with a finite affine critical set}
\label{sec:ordinaryfinite}

The three ordinary-node families admit one argument whenever the
critical set of $P$ is finite. In particular we need not solve for
each escaping critical point, assume it is Morse, or separate its
action from the actions of other escaping points. The Gaussian
quadratic itself detects the unwanted contribution at infinity.

Retain the notation and the full pair $(P,q)$ of
Lemma~\ref{lem:ordinarynormal}. Write $q_2$ for the homogeneous
quadratic part of $q$. The absence of a Gaussian frame is exactly
the following condition at the leading nodes:
\begin{equation}\label{eq:nodenoframe}
 F_2(w)=0\quad\Longrightarrow\quad q_2(w)\ne0.
\end{equation}
Indeed the Gaussian bilinear form in these affine coordinates is
$2q_2$; the value of $F_2$ at a leading singular direction is
unchanged by translation. In the common normal form the $T$ node
already satisfies this condition because $m\ne0$. The only extra
conditions are $c\ne0$ or $B\ne0$ when $d=0$, and $b\ne0$ or
$A\ne0$ when $a=0$.

\begin{lemma}[The stationary action identity]\label{lem:actionidentity}
Put $r=1/t$, and let $x(r)$ be an algebraic critical branch of
$S=P/r-q$. A prime means total differentiation along the branch.
Then
\begin{equation}\label{eq:actionidentity}
 S'=-P(x(r))/r^2,\qquad q(x(r))=-S-rS'.
\end{equation}
If $x(r)$ is bounded, the polar part of $S$ is $\alpha/r$, where
$\alpha=P(x(0))$. If it escapes and \eqref{eq:nodenoframe} holds,
its polar part is neither zero nor $\alpha/r$ for any constant
$\alpha$.
\end{lemma}
\begin{proof}
Use a finite Puiseux cover so every branch is meromorphic. At a
critical point the chain-rule term $\nabla_x S\cdot x'$ vanishes,
giving both identities. For a monomial $c r^{-\kappa}$ of the
polar part of $S$, its contribution to $-S-rS'$ is
$(\kappa-1)c r^{-\kappa}$. Thus bounded $q(x(r))$ permits only
the pole $\alpha/r$ in $S$. Its coefficient is $P(x(0))$ when
$x$ is bounded.

For an escaping branch let $R=\|x\|$. The critical equation is
$\nabla P=r\nabla q$, so $\|\nabla P\|/R\longrightarrow0$.
The leading projective direction is a singular point of $F_3$.
The ordinary-node argument in the proof of the preceding coercive
lemma works at this one direction: write $x=sw+y$, $y=o(s)$,
in a fixed complementary plane. The injectivity there of
$H=\operatorname{Hess}F_3(w)$ forces $y$ bounded from
\[
 \nabla P/s=Hy+\nabla F_2(w)+O((1+\|y\|)^2/|s|)=o(1).
\]
Pairing with $w$ then gives $F_2(w)=0$. Condition
\eqref{eq:nodenoframe} gives $q_2(w)\ne0$, so $q(x(r))$ has a
pole: its leading quadratic term cannot be cancelled by its
lower-degree terms. Identity~\eqref{eq:actionidentity} now rules
out both claimed polar parts. This argument includes all escape
scales and all coincident actions.
\end{proof}

We specify the classical rank input before using it. This is a
convergent algebraic application of formal stationary phase, not
an extension to arbitrary nonconvergent coefficients or rings.

\begin{lemma}[Rank accounting for a tame algebraic pencil]
\label{lem:tamepencilrank}
For the present no-frame pencil, shrink a punctured $r$ disk to
avoid its finitely many exceptional parameters. Its top twisted
de Rham module has rank
\[
 N=\sum_{\text{critical points of }S_r}\mu.
\]
After a finite cover, the rank of its block with nonzero
exponential polar part $\gamma$ is the sum of these generic local
Milnor numbers over critical branches with action polar part
$\gamma$. Equal actions add their multiplicities. Consequently
the rank of the zero-exponential block is the sum of the generic
Milnor numbers of the branches with bounded action.
\end{lemma}
\begin{proof}
For fixed small nonzero $r$, the homogeneous quadratic term of
$S_r$ at a leading node $w$ is $F_2(w)/r-q_2(w)$. It is nonzero
at every node, by \eqref{eq:nodenoframe} and exclusion of finitely
many parameters. The ordinary-node coercivity lemma gives
$\|\nabla S_r\|\ge c_r\|x\|$ outside a compact set. The
critical locus is therefore finite: a positive-dimensional affine
algebraic component would be unbounded. Pham's tame Milnor
decomposition~\cite[Part~I, Sections~1--4]{Pham} gives the displayed
rank, retaining the Milnor number of every degenerate point.

Here are the checks needed to apply the direct-image rank formula.
Take $U=\C_r^*\times\C_x^3$, $f=r$, $g=S_r$, and the regular
algebraic module $\mathcal O_U$. Form the ordinary direct image
by $(f,g)$, followed by exponential integration in the $g$
coordinate. All these are algebraic maps; $g$ is regular on $U$.
The ordinary direct image is regular holonomic. Over the punctured
$r$ disk its relevant cohomology is the middle, reduced
Gauss--Manin module of the fibers of $S_r$. The other, constant
fiber cohomology has zero exponential integral in the affine $g$
line. Punctual parameter terms disappear on the punctured disk.

There are no extra finite-$g$ singular curves from spatial infinity.
The coercive estimate is locally uniform on a compact subset of
the punctured $r$ disk. For $g$ in a compact set, the identity
$P=r(g+q)$ on a fiber gives
\[
 \partial_r S_r=-(g+q)/r=O(1+\|x\|^2).
\]
Lifting $r$ variation by the normalized gradient therefore uses a
vector field of at most linear growth. Lifting $g$ variation uses
a field $O(1/\|x\|)$. These fields are complete outside a fixed
ball on each such parameter neighborhood. They trivialize the
outer fiber pairs, so all nonvertical finite-$g$ discriminant
curves come from affine critical values. Along such a curve the
conormal multiplicity of the middle Gauss--Manin module is the
sum of the local Milnor numbers: a transverse target disk computes
precisely the corresponding isolated vanishing cycles. This also
applies to persistent degenerate points and equal critical values.

First make a common finite Puiseux base change of all the
critical-value curves. Apply the nonramified case of Roucairol's
Theorem~0.1(2), proved in her Section~1.2, and Theorem~0.3
\cite{RoucairolFormal} to this identified ordinary direct image.
They give exactly the stated nonzero polar parts and ranks.
Count every graph after this base change, including distinct
conjugates with the same polar part. For example the two graphs
$g=r^{-1}\pm r^{1/2}$ must contribute twice after $r=\rho^2$;
retaining just their one polar expression would lose rank.
Their rank statement requires no separation of the constant terms
of equal polar parts. The additional separation hypothesis in her
Theorem~0.2 concerns its monodromy formula, which we do not use.
Vertical components at $r=0$ and the boundary $g=\infty$ do not
add polar branches in Theorem~0.1. Subtracting the nonzero-block
ranks from $N$ proves the last assertion. In particular we have
not inferred the regular rank from an Euler characteristic alone.
\end{proof}

\begin{proposition}\label{prop:ordinaryfinite}
Suppose the common ordinary-node pair satisfies
\eqref{eq:nodenoframe} and $P$ has finite affine critical set.
Then $P$ is not pure. More generally every original Gaussian
marked period is either zero or is not a finite sum of
polynomials in $t$ times exponentials $e^{\beta t}$.
\end{proposition}
\begin{proof}
Every bounded critical branch of $S_r$ limits to a critical point
of $P$. Disjoint small Milnor neighborhoods of these finitely many
points contain, for the deformation $P-rq$, total critical
multiplicity $\mu_p$ in the neighborhood of $p$. This is
conservation of the local intersection multiplicity of the
gradient, and does not require a Morse perturbation. By
Lemma~\ref{lem:actionidentity} its action has polar part
$P(p)/r$. Every escaping branch has nonzero polar part. Hence
Lemma~\ref{lem:tamepencilrank} identifies the entire zero-exponential
rank as
\begin{equation}\label{eq:ordinaryzerorank}
 \sum_{p:\,dP(p)=0,\,P(p)=0}\mu_p.
\end{equation}

We construct that many independent local formal functionals and
compute their monodromy over the unramified field $\C((r))$.
At an isolated critical point at least one of $u,v$ is nonzero.
Indeed if both vanish, criticality requires $h=e=f=0$, in which
case the whole $T$ axis is critical. In the chart $u\ne0$ put
\[
 s=uv+h,\quad R(u)=G(u,-h/u),\quad
 T'=T+\frac{G(u,(s-h)/u)-G(u,-h/u)}s,
\]
where $G=P-T(uv+h)$. The quotient is analytic, and exactly
$P=sT'+R(u)$. The coordinate Jacobian is $u\ne0$. At a critical
point $p$ the one-variable critical multiplicity of $R$ is
$\mu_p$. An analytic one-variable coordinate $y$, with
$w=T'$, therefore gives the exact identity
\[
 P=\alpha+sw+y^k,\qquad k=\mu_p+1,\quad\alpha=P(p).
\]
The volume Jacobian and $e^{-q}$ multiply every amplitude by an
analytic unit $V(y,s,w)$; no Gaussian mean or cross term has been
discarded.

For $j=1,\ldots,k-1$ define a formal local functional on an
analytic monomial $y^a s^b w^c$ as zero unless $b=c$ and
$a=j-1+\ell k$ for an integer $\ell\ge0$. Otherwise its value is
\begin{equation}\label{eq:ordinarycharacters}
 e^{\alpha/r}r^{1+j/k}
       (-1)^{b+\ell}b!\,(j/k)_\ell\,r^{b+\ell}.
\end{equation}
Evaluate a polynomial form after multiplying its transformed
amplitude by $V$. At each power of $r$ only finitely many Taylor
terms contribute. The recurrences
\[
 I_{a+k}=-\frac{a+1}{k}r I_a,\qquad
 J_{b+1,b+1}=-(b+1)rJ_{b,b}
\]
prove spatial integration by parts and parameter horizontality.
Thus these are functionals on the full twisted quotient, not on
an assumed cyclic image. They exist even if the corresponding
asymptotic series diverge; their input coefficients here are the
Taylor series of the stated analytic coordinates.

They are independent. Prescribe, at each $p$, the jets
$1,y,\ldots,y^{k-2}$ on the line $s=w=0$ and zero jets at the
other points, using polynomial representatives of finite local
jets and the Chinese remainder theorem. After dividing column
$j$ by $e^{\alpha/r}r^{1+j/k}$, the leading evaluation matrix is
the jet map multiplied by the analytic unit $V$. It is triangular
with nonzero diagonal. Its higher terms are in $r\C[[r]]$.
For $\alpha=0$ these independent characters fill exactly the rank
in \eqref{eq:ordinaryzerorank}; there is no unaccounted regular
piece. Their formal monodromy eigenvalues, or their inverses if
one uses the dual module, are
\[
 \exp(2\pi i j/k),\qquad 1\le j<k.
\]
None is one. This conclusion is over $\C((r))$. We do not replace
a noninteger residue by its integer multiple after ramification.
The unique zero-exponential projector commutes with descent, and
the displayed unramified coefficient matrices identify its whole
block. Every other block has a nonzero exponential polar part.
The ambient module therefore has no trivial composition factor.

For completeness this excludes the formal Gaussian unit without
choosing an actual analytic period. Formal Gaussian integration
annihilates $\partial_i+tP_i-q_i$ coefficientwise. If its value
on a fixed marker $A$ were a nonzero constant, restriction to the
rational differential submodule generated by $A$ would give a
surjection to the trivial module: apply the functional to every
rational relation between $A,\nabla_tA,\ldots$. A relation has
zero constant coefficient because all positive derivatives of
that constant vanish. Exact formal extension at infinity would
preserve this composition factor, contradicting what was proved.

Finally a nonzero finite exponential-polynomial can be reduced
to $c e^{\beta t}$ by a constant-coefficient differential
operator. On an original marker this operator replaces $A$ by
the fixed polynomial marker $T(P)A$. Replace $P$ by $P-\beta$
and divide the marker by $c$. The hypotheses and finite critical
set are unchanged, so the same contradiction applies. Taking
$A=1$ proves the purity exclusion.
\end{proof}

\subsection{Exactly which ordinary-node cases remain}

The finite-critical hypothesis in Proposition~\ref{prop:ordinaryfinite}
has a short exhaustive coefficient test. This prevents a generic
argument from silently discarding the positive-dimensional cases.

If $h\ne0$, the critical locus lies on $uv=-h$, and the equations
above identify it with $R'(u)=0$ in $\C^*$. The Laurent polynomial
\[
 R(u)=a u^3+b u^2+e u+k-fh/u+ch^2/u^2-dh^3/u^3
\]
is nonconstant unless $a=b=c=d=e=f=0$. A nonzero Laurent
derivative has only finitely many roots in $\C^*$. In the constant
case the critical set is the smooth conic $uv=-h$, $T=0$.

If $h=0$, outside the center the critical equations are
\[
 \begin{array}{lll}
 v=0,\ u\ne0:&3au^2+2bu+e=0,&T=-f/u,\\
 u=0,\ v\ne0:&3dv^2+2cv+f=0,&T=-e/v.
 \end{array}
\]
The center $u=v=0$ is critical exactly when $e=f=0$, and then
the whole $T$ axis is critical. Consequently the following table
is exhaustive and disjoint:
\begin{center}\small
\begin{tabular}{@{}p{0.51\textwidth}p{0.39\textwidth}@{}}
\toprule
Condition & Conclusion\\
\midrule
$h\ne0$, $(a,b,c,d,e,f)\ne0$ & Finite critical set; excluded above\\
$h\ne0$, $a=b=c=d=e=f=0$ & Critical conic; still to be proved here\\
$h=0$, $e=f=0$ & Critical center axis, possibly extra curves; open here\\
$h=0$, $(e,f)\ne0$, $a=b=e=0$ & Critical hyperbola $v=0$, $uT=-f$; open here\\
$h=0$, $(e,f)\ne0$, $d=c=f=0$ & Symmetric critical hyperbola; open here\\
$h=0$, $(e,f)\ne0$, $(a,b,e)\ne0$, $(d,c,f)\ne0$
 & Finite critical set; excluded above\\
\bottomrule
\end{tabular}
\end{center}
The two hyperbola rows cannot overlap because $(e,f)\ne0$.
All isolated degeneracies in the finite rows have already been
included through their Milnor multiplicities. The curve rows
still need a full regular-block or homology argument; a critical
point count or the action identity alone does not settle them.
\begin{thebibliography}{9}
\bibitem{Pham}F.~Pham,\emph{La descente des cols par les onglets de
Lefschetz,avec vues sur Gauss--Manin},\emph{Ast\'erisque}
\textbf{130}(1985),11--47,Part~I,Sections~1--4;
\href{https://numdam.org/item/AST_1985__130__11_0/}{primary text}.
\bibitem{RoucairolFormal}C.~Roucairol,\emph{Formal structure of direct
image of holonomic $\mathcal D$-modules of exponential type},
\emph{Manuscripta Mathematica}\textbf{124}(2007),299--318,
Theorems~0.1(2),0.3 and Section~1.2;
\href{https://arxiv.org/abs/math/0604134v2}{arXiv:math/0604134v2}.
\end{thebibliography}\end{document}
